% docs/texte/05_theory/notation.tex -- theory chapter "Notation" (Interferosim)
% Version: 0.1.1 (2026-10-12)
% 0.1.1: chapter renamed "Notation and Conventions" -> "Notation";
%        paragraph "Conventions" -> "General assumptions" (no confusion with
%        the coding conventions in part "Software Engineering")
% Notation as defined in texte/05_theory/notation.tex.
% ================================================================

\index{notation}
This part collects the physical and mathematical basis of
Interferosim. All chapters use the notation of
Table~\ref{tab:notation}; the same symbols are used in the code
documentation and in the notebooks. Standard references are
\cite{thompson2017,condon2016,wilson2013,taylor1999}.

\paragraph*{General assumptions.}
\begin{itemize}
  \item SI units throughout; flux densities in jansky,
        $1\,\mathrm{Jy} = 10^{-26}\,\mathrm{W\,m^{-2}\,Hz^{-1}}$;
        solar flux unit $1\,\mathrm{SFU} = 10^{4}\,\mathrm{Jy}$.
  \item Spatial frequencies $u,v,w$ are measured in \emph{wavelengths},
        baselines $\vec b$ in metres.
  \item Sky coordinates $(l,m)$ are direction cosines relative to the
        phase centre; for small fields they equal angular offsets in
        radians.
  \item The Fourier sign convention is $V = \int I\,e^{-2\pi i(ul+vm)}$
        (sky $\to$ visibility) and $I = \int V\,e^{+2\pi i(ul+vm)}$
        (visibility $\to$ sky).
  \item Brightness temperatures use the Rayleigh--Jeans approximation
        ($h\nu \ll kT$). At 1.4\,GHz, $h\nu/k \approx 0.07\,\mathrm{K}$;
        the relative error $\approx h\nu/2kT$ is below $10^{-3}$ for
        $T > 35\,\mathrm{K}$, which holds for all sources considered.
\end{itemize}

\begin{longtable}{@{}>{$}l<{$}l>{\raggedright\arraybackslash}p{8.2cm}@{}}
\caption{Symbols used in this documentation}\label{tab:notation}\\
\toprule
\textrm{Symbol} & Unit & Meaning \\
\midrule
\endfirsthead
\toprule
\textrm{Symbol} & Unit & Meaning \\
\midrule
\endhead
\bottomrule
\endfoot
\multicolumn{3}{@{}l}{\textit{Constants}}\\
c & m\,s$^{-1}$ & speed of light, $299\,792\,458$ \\
k & J\,K$^{-1}$ & Boltzmann constant, $1.380649\times10^{-23}$ \\
\omega_E & rad\,s$^{-1}$ & Earth rotation rate, sidereal
  $7.292115\times10^{-5}$; for the Sun $\omega_\odot = 2\pi/86\,400\,\mathrm{s}$ \\
\multicolumn{3}{@{}l}{\textit{Radiation and source}}\\
\nu,\ \lambda & Hz, m & observing frequency and wavelength, $\lambda = c/\nu$ \\
\Delta\nu & Hz & bandwidth \\
I(l,m) & Jy\,sr$^{-1}$ & sky brightness (specific intensity) \\
T_b & K & brightness temperature \\
S & Jy & flux density \\
\Omega & sr & solid angle (of a source) \\
\theta_d & rad & diameter of a uniform disc \\
\theta_{1/2} & rad & FWHM of a Gaussian source \\
\multicolumn{3}{@{}l}{\textit{Antenna and receiver}}\\
D & m & dish diameter \\
\eta_A & -- & aperture efficiency \\
A_{\mathrm{eff}} & m$^2$ & effective area \\
G & -- & gain (often quoted in dBi) \\
\theta_{\mathrm{HPBW}} & rad & half-power beam width \\
k_b & -- & beam factor, $\theta_{\mathrm{HPBW}} = k_b\,\lambda/D$ \\
A(l,m) & -- & normalised primary beam \\
T_A,\ T_{\mathrm{sys}} & K & antenna and system temperature \\
\mathrm{SEFD} & Jy & system equivalent flux density \\
\tau & s & integration time \\
\rho & -- & correlation coefficient \\
\multicolumn{3}{@{}l}{\textit{Geometry}}\\
\vec b,\ b & m & baseline vector and length \\
L_X, L_Y, L_Z & m & baseline components in equatorial coordinates \\
\hat s & -- & unit vector towards the source \\
\tau_g & s & geometric delay, $\tau_g = \vec b\cdot\hat s/c$ \\
\theta & rad & angle between baseline and source direction \\
H,\ \delta & rad & hour angle and declination of the source \\
\varphi & rad & geographic latitude of the site \\
\varepsilon & rad & elevation \\
\multicolumn{3}{@{}l}{\textit{Fourier domain and imaging}}\\
u,v,w & $\lambda$ & spatial frequencies (baseline in wavelengths) \\
q & $\lambda$ & radial spatial frequency, $q = \sqrt{u^2+v^2}$ \\
l,m & -- & direction cosines \\
\mathcal{V}(u,v) & Jy & complex visibility \\
\mathcal{S}(u,v) & -- & sampling (transfer) function \\
W(u,v) & -- & imaging weights \\
I_D,\ B_D & Jy\,beam$^{-1}$, -- & dirty image, dirty beam \\
B_C & -- & clean beam \\
\mathcal{W} & -- & Wiener filter \\
\Phi_s,\ \Phi_n & -- & power spectra of signal and noise \\
f_F & Hz & fringe frequency (fringe rate) \\
\end{longtable}
